\documentclass[11pt]{article}
\usepackage[T1]{fontenc}
\usepackage{lmodern}
\usepackage[margin=1in]{geometry}
\usepackage{amsmath,amssymb,amsthm}
\usepackage{booktabs}
\setlength{\emergencystretch}{2em}
\usepackage[colorlinks=true,linkcolor=blue,urlcolor=blue]{hyperref}
\newtheorem{definition}{Definition}
\newtheorem{theorem}{Theorem}
\newtheorem{proposition}{Proposition}
\newcommand{\Run}{\ensuremath{\mathsf{Run}}}
\newcommand{\Trace}{\ensuremath{\mathsf{TraceRun}}}
\newcommand{\Step}{\ensuremath{\mathsf{Step}}}
\newcommand{\erase}{\ensuremath{\mathsf{outputs}}}
\newcommand{\silent}{\ensuremath{\mathsf{silentCount}}}
\newcommand{\None}{\ensuremath{\mathsf{none}}}
\newcommand{\Some}{\ensuremath{\mathsf{some}}}
\newcommand{\cat}{\mathbin{+\!+}}
\title{Preserving Silent Calls in SPQR Runs\\A Two-Agent Definition and Review Experiment}
\author{Local research experiment\\AI-assisted development and independent agent review}
\date{10 November 2026}
\begin{document}
\maketitle
\begin{abstract}
The existing SPQR verification development represents a finite run by its initial state,
final state, and list of emitted epoch secrets. Successful calls that emit no key contribute
no entry to that list. We introduce a conservative refinement recording one optional output
per successful call. We prove that existential erasure of these traces is exactly the existing
run relation, that every specified trace boundary admits an intermediate state, and that
trace length equals emitted-key count plus silent-call count. The definitions and proofs are
checked in Lean against the repository's existing transition relation. A second agent reviews
the construction, its trust dependencies, and the limits of its interpretation. This is a
small local formalization experiment, with no claim of literature novelty, optimality, or
additional cryptographic security.
\end{abstract}

\section{Problem and experimental method}
A useful definition exposes the information needed by later proofs without changing the
behaviors already modeled. Here the target is the relation \Run\ in the local SPQR
verification repository~\cite{run}. Its output index records emitted keys in order, but has
no position for a successful call producing no key. Induction on a derivation can still
reason about such calls; the issue is what information the relation's explicit indices
expose, not the absence of silent transitions from the existing semantics.

The user requested two agents: one to improve or introduce definitions and derive theorems,
and another to review those results. The author agent produced a separate Lean file and
an author report. The reviewer examined the existing semantics, checked the new statements
and proofs, and produced a separate review. The coordinating agent assembled this paper and
its reproducible artifacts. These roles provide a review workflow, not evidence of statistical
independence or a substitute for human review.

We interpret ``better'' relative to two concrete criteria: retaining successful-call positions
and preserving exactly the old relation after erasure. We do not optimize over all possible
definitions. Production definitions remain unchanged. All experimental files live under
\texttt{experiments/two-agent-traces/}, on branch
\path{experiment/two-agent-trace-definitions}.

\section{Existing semantics and the refined definition}
Let $S$ be the existing type \texttt{States} and $K$ the type \texttt{EpochSecret}.
Write $\Step(s,t,o)$ for the existing relation with $o\in\operatorname{Option}(K)$.
It holds when a send or receive call succeeds at both the Aeneas and Rust result levels,
returns state $t$, and returns optional key $o$. Send witnesses include RNG interfaces and
states; receive witnesses include the message. These witnesses are not explicit indices of
our trace either. The relation concerns one party, arbitrary initial states, and successful
calls only.

For an optional key $o$, put $\operatorname{list}(\None)=[]$ and
$\operatorname{list}(\Some(k))=[k]$. The existing relation has constructors
\[
\frac{}{\Run(s,s,[])}\qquad
\frac{\Step(s,u,o)\quad\Run(u,t,ks)}
 {\Run(s,t,\operatorname{list}(o)\cat ks)}.
\]
Consequently, a silent call contributes no element to the emitted-key list.

\begin{definition}[Successful-call trace]
A trace is a list $\tau\in\operatorname{List}(\operatorname{Option}(K))$.
Define the relation $\Trace$ by
\[
\frac{}{\Trace(s,s,[])}\qquad
\frac{\Step(s,u,o)\quad\Trace(u,t,\tau)}
 {\Trace(s,t,o::\tau)}.
\]
Define erasure recursively by
$\erase([])=[]$ and
$\erase(o::\tau)=\operatorname{list}(o)\cat\erase(\tau)$.
\end{definition}

Every trace entry now corresponds to one successful call, including a silent call.
The explicit trace does not distinguish send from receive, record messages, or constrain a
sequence of RNG states across calls beyond what the original relation requires. It is an
optional-output history, not a full execution log. No uniqueness of intermediate states is asserted.

\section{Conservativity and composition}
\begin{theorem}[Exact recovery of the existing relation]
For all $s,t\in S$ and $ks\in\operatorname{List}(K)$,
\[
\Run(s,t,ks)\quad\Longleftrightarrow\quad
\exists\tau,\ \Trace(s,t,\tau)\ \land\ \erase(\tau)=ks.
\]
\end{theorem}
\begin{proof}
For the forward direction, induct on the old run. The empty run lifts to the empty trace.
For a constructor with head output $o$, lift the tail to $\tau$ and prepend $o$.
The erasure equation follows from its recursive definition. For the reverse direction,
induct on the trace derivation: each constructor becomes the corresponding old run
constructor, whose index prepends $\operatorname{list}(o)$ rather than $o$.
Substitute the supplied erasure equality.
\end{proof}

This is a logical equivalence of relations, with an existential trace witness. It does not
assert a unique trace or a canonical executable extraction function from proof terms.
It establishes compatibility with every property already stated purely over the old run's
endpoints and emitted keys: given such a property for \Run, apply it to the erased trace.
This observation does not reprove any existing protocol invariant.

\begin{proposition}[Erasure and composition]
For all traces $x,y$,
$\erase(x\cat y)=\erase(x)\cat\erase(y)$.
If $\Trace(s,u,x)$ and $\Trace(u,t,y)$, then $\Trace(s,t,x\cat y)$.
\end{proposition}
\begin{proof}
The erasure equation follows by induction on $x$ and associativity of list concatenation.
The relational claim follows by induction on the first trace derivation, retaining each
head transition and appending the second derivation at its empty tail.
\end{proof}

\section{Exact decomposition at call boundaries}
\begin{theorem}[Trace splitting]
For all states $s,t$ and traces $x,y$,
\[
\Trace(s,t,x\cat y)\quad\Longleftrightarrow\quad
\exists u,\ \Trace(s,u,x)\land\Trace(u,t,y).
\]
\end{theorem}
\begin{proof}
For the forward direction induct on $x$, generalizing $s$. If $x=[]$, take $u=s$.
If $x=o::x'$, invert the trace constructor to obtain its first transition and remaining
trace. Apply the induction hypothesis to that remainder and prepend the first transition
to its prefix. The reverse direction is trace composition.
\end{proof}

\begin{theorem}[Cut after a specified number of calls]
If $\Trace(s,t,\tau)$, then for every natural number $n$ there is a state $u$ such that
\[
\Trace(s,u,\operatorname{take}(n,\tau))\ \land\
\Trace(u,t,\operatorname{drop}(n,\tau)).
\]
\end{theorem}
\begin{proof}
Use $\operatorname{take}(n,\tau)\cat\operatorname{drop}(n,\tau)=\tau$
and apply trace splitting.
\end{proof}

For $n\leq |\tau|$, the prefix contains exactly $n$ successful calls; for $n>|\tau|$,
list operations saturate, giving the whole trace and an empty suffix. The theorem provides
existence, not uniqueness, of $u$. A boundary counts successful calls, not elapsed time,
network deliveries, failed calls, or emitted keys.

As a list illustration, take
$\tau=[\None,\Some(k_1),\None,\Some(k_2)]$.
Its erasure is $[k_1,k_2]$, and a cut at three calls retains
$[\None,\Some(k_1),\None]$ as the prefix. If this list has a trace derivation, the cut
theorem supplies a corresponding intermediate state. The example itself does not establish
that this particular list is realizable by SPQR.

\section{Exact accounting and information loss}
\begin{definition}[Silent-call count]
Define $\silent([])=0$, $\silent(\None::\tau)=1+\silent(\tau)$, and
$\silent(\Some(k)::\tau)=\silent(\tau)$.
\end{definition}
\begin{theorem}[Call accounting]
For every optional-output list $\tau$,
\[
|\tau|=|\erase(\tau)|+\silent(\tau),\qquad |\erase(\tau)|\leq|\tau|.
\]
\end{theorem}
\begin{proof}
Induct on $\tau$ and split on its head option. A \None\ increments only the silent count;
a \Some\ increments only the emitted-key count. Natural-number arithmetic yields the
equality, and nonnegativity yields the inequality.
\end{proof}
For a valid trace, this list identity becomes an exact accounting statement about successful
calls. It is not a progress bound: neither the definition nor this theorem bounds the number
of silent calls needed to produce a key.

Erasure is not injective on unrestricted lists, since $\erase([])=\erase([\None])=[]$
although the two lists differ. This is separately proved in Lean. It does not prove the
existence of two realizable traces with the same endpoints and erased outputs. In particular,
no silent self-loop is assumed. This distinction is a central constraint from the review.

\section{Formal artifact and review}
All new declarations are in the namespace
\path{spqr.v1.chunked.states.TraceExperiment}. Table~\ref{tab:map} maps mathematical
results to their Lean declarations; names in the table are relative to that namespace.
\begin{table}[ht]
\centering\small
\begin{tabular}{ll}
\toprule
Result & Lean declaration \\
\midrule
Conservativity & \texttt{run\_iff\_exists\_trace} \\
Erasure of a valid trace & \texttt{TraceRun.toRun} \\
Lifting an old run & \texttt{Run.exists\_trace} \\
Erasure of concatenation & \texttt{outputs\_append} \\
Composition & \texttt{TraceRun.append} \\
Splitting equivalence & \texttt{traceRun\_append\_iff} \\
Arbitrary call cut & \texttt{TraceRun.cut} \\
Exact call count & \texttt{length\_eq\_outputs\_add\_silent} \\
Lower bound on call count & \texttt{outputs\_length\_le} \\
List erasure noninjectivity & \texttt{outputs\_not\_injective} \\
\bottomrule
\end{tabular}
\caption{Correspondence between the paper and checked declarations.}\label{tab:map}
\end{table}

The author supplies \texttt{TraceRun.lean} and \texttt{AUTHOR.md}; the reviewer supplies
\texttt{REVIEW.md} and separate checks. The review assesses conservativity, successful-call
scope, existential rather than unique witnesses, and the difference between arbitrary lists
and executable histories. Final commands and trust-audit results are recorded with the
artifact so that validation is reproducible rather than inferred from prose.

The pinned repository toolchain is Lean \texttt{v4.35.0-rc4}.\par
Base revision: {\small\texttt{e38c4ce87a10c12d709a79962d75ac03ec295182}}.\par
The experiment imports the existing run module and adds no new axioms or admitted proofs.
That statement concerns the new development: it does not certify all external cryptographic
functions or every assumption elsewhere in the repository. The independent axiom audit reports \texttt{propext}, \texttt{Classical.choice},
\texttt{Quot.sound}, \path{Aeneas.Std.core.fmt.Formatter}, and 19 inherited
private translated \texttt{\_native.decide.ax\_N} axiom constants for the seven run-related
results. No audited result depends on \texttt{sorryAx}. Pure list results have smaller
dependency sets, detailed in the review. Thus these proofs are not axiom-free.
The mathematical proofs use
the existing step relation parametrically, without proving that a concrete send or receive
call succeeds.

\section{Limitations and outcome}
This experiment establishes a small refinement suited to reasoning at successful-call
boundaries. The old definition remains useful for properties depending only on emitted
keys. The refinement costs an additional trace index and supporting lemmas, and does not
supply liveness, computational security, two-party agreement, or resource bounds. It does
not record failed calls or explicit action labels. A richer action trace would require a
separate definition and another conservativity argument.

The new results are additions to this experiment, not claims of mathematical novelty in the
literature. They use elementary inductive arguments, now checked against the concrete
repository relation. The review narrows their interpretation and documents what would
require further evidence. No global ``best definition'' claim follows from the experiment.

\begingroup
\small
\raggedright
\begin{thebibliography}{9}
\bibitem{run} SPQR verification repository, local base revision cited above.
\texttt{Spqr/Specs/Braid/StateMachine/Run.lean}.
Existing definitions of \texttt{Step}, \texttt{Run}, and run composition.
\bibitem{artifact} Two-agent trace experiment, accompanying local artifacts.
\texttt{experiments/two-agent-traces/TraceRun.lean},
\texttt{AUTHOR.md}, and \texttt{REVIEW.md}, 2026.
\end{thebibliography}
\endgroup
\end{document}
